% !TeX program = xelatex
% 编译：xelatex 量子力学入门讲义.tex（运行两次以生成目录和交叉引用）
% 使用常规 TeX Live / MiKTeX 中的 xeCJK、amsmath 等宏包，无外部图片或参考文献文件。
\documentclass[a4paper,11pt]{article}
\usepackage{fontspec}
\usepackage{xeCJK}
\IfFontExistsTF{FandolSong-Regular}{
  \setCJKmainfont{FandolSong-Regular}[BoldFont=FandolSong-Bold]
}{
  \setCJKmainfont{SimSun}[AutoFakeBold=2.2]
}
\setmainfont{Latin Modern Roman}
\usepackage[margin=25mm,top=23mm,bottom=24mm]{geometry}
\usepackage{amsmath,amssymb,mathtools}
\usepackage{enumitem,booktabs,array}
\usepackage{titlesec,needspace}
\usepackage[unicode,hidelinks]{hyperref}
\hypersetup{pdftitle={量子力学入门：从线性空间到时间演化},pdfauthor={}}
\renewcommand{\contentsname}{目录}
\renewcommand{\abstractname}{阅读说明}
\renewcommand{\figurename}{图}
\renewcommand{\tablename}{表}
\titleformat{\section}{\large\bfseries}{\thesection}{0.7em}{}
\titleformat{\subsection}{\normalsize\bfseries}{\thesubsection}{0.7em}{}
\titlespacing*{\section}{0pt}{2.0ex plus .5ex minus .2ex}{1.0ex}
\titlespacing*{\subsection}{0pt}{1.5ex plus .3ex minus .2ex}{.6ex}
\setlength{\parindent}{2em}
\setlength{\parskip}{.22em}
\linespread{1.13}
\setlist{nosep,topsep=.5em,leftmargin=2em}
\setcounter{tocdepth}{1}
\numberwithin{equation}{section}
\allowdisplaybreaks[1]
\raggedbottom
\emergencystretch=1.5em
\newcommand{\ii}{\mathrm{i}}
\newcommand{\dd}{\mathop{}\!\mathrm{d}}
\newcommand{\ee}{\mathrm{e}}
\newcommand{\HH}{\mathcal H}
\newcommand{\I}{\mathbb I}
\newcommand{\ket}[1]{\lvert #1\rangle}
\newcommand{\bra}[1]{\langle #1\rvert}
\newcommand{\braket}[2]{\langle #1\mid #2\rangle}
\newcommand{\mel}[3]{\langle #1\rvert #2\lvert #3\rangle}
\newcommand{\avg}[1]{\langle #1\rangle}
\newcommand{\comm}[2]{[#1,#2]}
\newcommand{\norm}[1]{\lVert #1\rVert}
\newcommand{\abs}[1]{\lvert #1\rvert}
\newcommand{\important}[1]{\par\noindent\textbf{#1}\quad}
\title{\vspace{-1.5em}\textbf{量子力学入门}\\[.35em]\large 从线性空间到时间演化}
\author{}
\date{}

\begin{document}
\maketitle
\vspace{-2.5em}
\begin{abstract}
本文面向学过线性代数与微积分、尚未接触量子力学的读者。旨在不引入大量物理背景的情况下，尽可能以最小的篇幅介绍量子力学的核心原理
\end{abstract}

\vspace{.6em}

\section{先回顾线性空间：向量不等于一列数字}
\subsection{向量、线性组合与基}
线性空间的核心，是其中的对象可以相加、可以乘以标量，而且这些运算满足熟悉的分配律等规则。对象可以是几何箭头、列向量，也可以是函数。例如，对两个复值函数 $f,g$，
\[
  (af+bg)(x)=af(x)+bg(x),\qquad a,b\in\mathbb C
\]
定义了它们的线性组合。量子力学中的态空间通常是\textbf{复线性空间}，因此不能只允许实数系数。

若 $e_1,\ldots,e_N$ 构成一组基，那么每个向量都唯一写成
\[
  v=\sum_{n=1}^N c_n e_n.
\]
列向量 $(c_1,\ldots,c_N)^{\mathsf T}$ 是 $v$ 在这组基下的\textbf{坐标}。换一组基，坐标就会改变，但抽象向量仍是同一个。以后所谓“位置表示”“动量表示”，本质上就是这一区别的延伸。

线性代数还研究线性映射 $A$：
\[
  A(av+bw)=aAv+bAw.
\]
当它把一个空间映到自身时，物理中常称它为\textbf{线性算符}。选定基以后，算符由矩阵表示；矩阵是算符在某组坐标中的写法。

\subsection{复内积：长度、正交性与展开系数}
为了描述长度和正交性，需要给线性空间配备内积。全文采用物理中的惯例：内积对\textbf{第一个变量共轭线性、第二个变量线性}。在 $\mathbb C^N$ 中，
\[
  (v,w)=\sum_{n=1}^N v_n^*w_n,\qquad
  (v,w)=(w,v)^*,\qquad \norm{v}^2=(v,v).
\]
星号表示复共轭。特别地，$(av,w)=a^*(v,w)$，而 $(v,aw)=a(v,w)$。若漏掉复共轭，连 $\norm{v}^2\geq0$ 都无法保证。

若基是正交归一的，即 $(e_m,e_n)=\delta_{mn}$，则
\[
  c_n=(e_n,v),\qquad \norm{v}^2=\sum_n\abs{c_n}^2.
\]
这两式将分别变成量子力学中的\textbf{概率幅}与\textbf{概率归一化}。

函数也可以配备内积，例如
\begin{equation}
  (f,g)=\int_{\mathbb R}f(x)^*g(x)\dd x,
  \qquad \norm{f}^2=\int_{\mathbb R}\abs{f(x)}^2\dd x.
\end{equation}
满足 $\norm{f}<\infty$ 的平方可积函数组成 $L^2(\mathbb R)$；严格说，几乎处处相等的函数在这里视为同一个向量。这个空间通常是无限维的。

\textbf{Hilbert 空间}就是完备的内积空间。“完备”保证按这个范数得到的柯西序列有空间内的极限。有限维复内积空间自动完备；$L^2$ 也是完备的。入门时，只需把它理解为允许我们可靠地做正交展开、取极限和计算内积的态空间。

\section{状态假设与 Dirac 符号}
\subsection{纯态、归一化与相位}
\important{状态假设}
一个量子系统对应一个复 Hilbert 空间 $\HH$。系统的\textbf{纯态}用一个非零向量所张成的射线表示；计算时通常取其中归一化的代表，记作 $\ket{\psi}$，满足
\begin{equation}
  \norm{\ket{\psi}}^2=1.
\end{equation}
$\ket{\psi}$ 读作“ket psi”，只是给抽象向量换了一种记号。

为什么说是“射线”？因为 $\ket{\psi}$ 与 $c\ket{\psi}$（$c\neq0$）归一化后只差整体相位，表示同一个纯态。对归一化向量，
\begin{equation}
  \ket{\psi}\sim\ee^{\ii\alpha}\ket{\psi},\qquad \alpha\in\mathbb R.
\end{equation}
这个相位乘在\textbf{整个向量}前面。不同分量之间的相位差则通常能改变测量概率。

例如，在二维态空间中，选取正交归一基 $\ket{0},\ket{1}$，一般的归一化态为
\begin{equation}\label{eq:qubit}
  \ket{\psi}=a\ket{0}+b\ket{1},\qquad \abs{a}^2+\abs{b}^2=1.
\end{equation}
$0,1$ 此时只是两种基态的标签，并不自动表示位置、能量或其他物理量。空间的线性结构允许叠加态；究竟怎样测量这个态，需要下面的测量假设。

\subsection{对偶空间与 bra：把向量变成一个数的规则}
对偶空间中的对象是\textbf{线性泛函}，即把向量线性地映到复数的函数，也就是有
\begin{equation}
    T(\ket{\psi}+\ket{\phi})=T\ket{\psi}+T\ket{\phi}
\end{equation}
例如，在有限维坐标中，固定一行数字 $r$ 后，$v\mapsto rv$ 就是线性泛函。线性泛函自己也组成了一个线性空间，叫做\textbf{对偶空间}

Riesz 表示定理保证：Hilbert 空间中的每个连续线性泛函都等价于和Hilbert空间中的某个元素作内积诱导出的泛函。这意味着我们可以把所有满足条件的泛函唯一写成 $\bra{\phi}$ 的形式。
\begin{equation}
  \bra{\phi}:\ket{\psi}\longmapsto\braket{\phi}{\psi}.
\end{equation}
$\bra{\phi}$ 读作“bra phi”，而 $\braket{\phi}{\psi}$ 就是两向量的内积。在正交归一基下，
\begin{equation}
  \ket{\phi}\ \longleftrightarrow\
  \begin{pmatrix}\phi_1\\ \vdots\\ \phi_N\end{pmatrix},
  \qquad
  \bra{\phi}\ \longleftrightarrow\
  \begin{pmatrix}\phi_1^*&\cdots&\phi_N^*\end{pmatrix}.
\end{equation}
因此 ket 变成 bra 时必须复共轭：
\begin{equation}
  a\ket{\phi}+b\ket{\psi}
  \quad\longmapsto\quad
  a^*\bra{\phi}+b^*\bra{\psi}.
\end{equation}
bra 对它所作用的 ket 是线性的；\textbf{从 ket 到对应 bra 的映射}却是共轭线性的。这是两个不同的陈述。


\subsection{正交展开、外积与单位算符}
若 $\{\ket{n}\}$ 是一组完备的正交归一基，那么
\begin{equation}\label{eq:expansion}
  \ket{\psi}=\sum_n c_n\ket{n},\qquad
  c_n=\braket{n}{\psi},\qquad \sum_n\abs{c_n}^2=1.
\end{equation}
无限和在 Hilbert 空间范数下理解。“完备”意味着没有遗漏任何独立的方向。

\textbf{“算符”}指的是Hilbert空间到他自己的映射，比如有限维线性空间$\mathbb{C}^n$上的$n\times n$矩阵也可以看成一个算符。实际上，我们在量子力学中会遇到很多可以写成矩阵的算符，因为很多关心的系统的状态本质上只在一个有限维的子空间里。

与内积 $\braket{\phi}{\psi}$ 不同，\textbf{外积} $\ket{\phi}\bra{\chi}$ 是一个算符。作用规则是
\begin{equation}
  (\ket{\phi}\bra{\chi})\ket{\psi}
  =\ket{\phi}\braket{\chi}{\psi}.
\end{equation}
先用 bra 取出一个复数，再把它乘到 ket 上。如果 $\ket{n}$ 已归一化，则
\[
  P_n=\ket{n}\bra{n}
\]
把任意向量投影到 $\ket{n}$ 张成的一维子空间。它满足 $P_n^2=P_n$：投影一次之后，再投影不会改变结果。

把所有方向上的投影相加，得到\textbf{完备性关系}：
\begin{equation}\label{eq:resolution}
  \boxed{\sum_n\ket{n}\bra{n}=\I.}
\end{equation}
$\I$ 是单位算符。这不是新的物理假设，而是“基完备”的另一种写法。以后需要切换表示或插入中间态时，最常用的技巧就是在适当位置插入这个单位算符。

\section{可观测量与测量：从算符得到概率}
\subsection{算符、矩阵元与伴随}
一个线性算符 $A$ 作用于 ket 后仍给出 ket：$A\ket{\psi}$。在正交归一基 $\{\ket{n}\}$ 下定义
\begin{equation}
  A_{mn}=\mel{m}{A}{n},\qquad
  A=\sum_{m,n}A_{mn}\ket{m}\bra{n}.
\end{equation}
所以 $A\ket{\psi}$ 的第 $m$ 个坐标为 $\sum_n A_{mn}c_n$，正是矩阵乘列向量。两个算符相乘时，右边先作用：$AB\ket{\psi}=A(B\ket{\psi})$。

\textbf{伴随算符} $A^\dagger$ 的定义是
\begin{equation}\label{eq:adjoint}
  \braket{\phi}{A\psi}=\braket{A^\dagger\phi}{\psi},
  \qquad\text{也即}\qquad
  (A\ket{\phi})^\dagger=\bra{\phi}A^\dagger.
\end{equation}
在正交归一基中，它对应共轭转置矩阵：$(A^\dagger)_{mn}=A_{nm}^*$。直接由定义可得
\begin{equation}
  (AB)^\dagger=B^\dagger A^\dagger,\qquad
  (\lambda A)^\dagger=\lambda^*A^\dagger,
  \qquad (\ket{\phi}\bra{\chi})^\dagger=\ket{\chi}\bra{\phi}.
\end{equation}
顺序倒转来自矩阵乘法，不能漏掉。

\important{可观测量假设}
每个可观测量由态空间中的一个\textbf{自伴算符}表示。在有限维中，自伴就是 $A^\dagger=A$，也称 Hermitian 或厄米；无限维时还必须比较算符及其伴随的定义域，不能只检查形式上的共轭转置。

这里的物理动机是：测量结果是实数，而自伴算符具有实谱，并能用正交投影分解。以本征向量为例，若 $A\ket{a}=a\ket{a}$ 且 $\braket{a}{a}=1$，则
\[
  a=\mel{a}{A}{a}=\mel{a}{A}{a}^*=a^*.
\]
若 $A\ket{a}=a\ket{a}$、$A\ket{b}=b\ket{b}$ 且 $a\neq b$，则
\[
  a\braket{a}{b}=\mel{a}{A}{b}=b\braket{a}{b},
  \qquad\Rightarrow\qquad\braket{a}{b}=0.
\]
不同本征值对应的本征态正交。

\subsection{测量结果、Born 规则与测量后的态}
先考虑离散谱。为容纳简并，用 $a$ 标记\textbf{不同的本征值}，用 $P_a$ 表示对应整个本征子空间的正交投影：
\begin{equation}\label{eq:spectral}
  A=\sum_a aP_a,\qquad
  P_aP_b=\delta_{ab}P_a,\qquad \sum_aP_a=\I.
\end{equation}
若本征值不简并，$P_a=\ket{a}\bra{a}$；若有简并，$P_a=\sum_r\ket{a,r}\bra{a,r}$。

\important{测量假设（理想投影测量）}
对归一化态 $\ket{\psi}$ 测量 $A$，可能结果是 $A$ 的谱值；在上述离散情形中，得到 $a$ 的概率为
\begin{equation}\label{eq:born}
  \boxed{p(a)=\mel{\psi}{P_a}{\psi}=\norm{P_a\ket{\psi}}^2.}
\end{equation}
若记录到结果 $a$ 且 $p(a)>0$，在只分辨这个本征值的理想 Lüders 测量中，条件态变为
\begin{equation}\label{eq:projection}
  \ket{\psi}\longmapsto
  \frac{P_a\ket{\psi}}{\sqrt{p(a)}}.
\end{equation}
因此立即重复同一测量会得到同一结果。归一化分母来自投影后向量的长度，而不是任意加入的修正。

非简并时，式\eqref{eq:born}化为熟悉的\textbf{Born 规则}
\begin{equation}
  p(a)=\abs{\braket{a}{\psi}}^2.
\end{equation}
$\braket{a}{\psi}$ 是概率幅，模平方才是概率。式\eqref{eq:projection}描述特定的理想测量模型；一般实际测量可以更复杂。尤其有简并时，测量后状态如何变化还取决于仪器是否进一步分辨子空间内的信息。

\subsection{平均值、涨落与确定值}
概率分布给出的平均测量结果为
\begin{equation}\label{eq:expectation}
  \avg{A}_\psi=\sum_a a\,p(a)=\mel{\psi}{A}{\psi}.
\end{equation}
它指大量\textbf{同样制备的系统}各测一次所得的统计平均，比如结果只有 $+1,-1$ 时，平均值可以是 $0$。

用标准差刻画分布宽度：
\begin{equation}\label{eq:variance}
  (\Delta A)^2=\avg{(A-\avg{A}\I)^2}
  =\avg{A^2}-\avg{A}^2
  =\norm{(A-\avg{A}\I)\ket{\psi}}^2.
\end{equation}
所以 $\Delta A=0$ 当且仅当该态属于一个确定本征值的本征子空间。一个系统能处于 $A$ 的确定值态，但不一定同时处于另一个可观测量的确定值态。

\subsection{一个例子看清叠加、相对相位与混合}
在二维空间定义两个无量纲可观测量
\begin{equation}\label{eq:ZX}
  Z=\ket{0}\bra{0}-\ket{1}\bra{1},\qquad
  X=\ket{0}\bra{1}+\ket{1}\bra{0}.
\end{equation}
$Z$ 的本征态为 $\ket{0},\ket{1}$；$X$ 的本征态则是
\begin{equation}\label{eq:pm}
  \ket{+}=\frac{\ket{0}+\ket{1}}{\sqrt2},\qquad
  \ket{-}=\frac{\ket{0}-\ket{1}}{\sqrt2},\qquad
  X\ket{\pm}=\pm\ket{\pm}.
\end{equation}
对于式\eqref{eq:qubit}，测量 $Z$ 得到 $+1,-1$ 的概率是 $\abs{a}^2,\abs{b}^2$；测量 $X$ 得到 $+1$ 的概率却是
\begin{equation}
  p_X(+1)=\abs{\braket{+}{\psi}}^2
  =\frac12\abs{a+b}^2
  =\frac12+\operatorname{Re}(a^*b).
\end{equation}
最后一项是干涉项，依赖 $a,b$ 的相对相位。

更具体地，令 $\ket{\psi_\varphi}=(\ket{0}+\ee^{\ii\varphi}\ket{1})/\sqrt2$。测量 $Z$ 始终得到各 $1/2$ 的概率，但
\[
  p_X(+1)=\frac{1+\cos\varphi}{2}.
\]
整体相位在模平方中消失，相对相位却改变了另一种测量的统计结果。

还要区分这个纯态与“以各 $1/2$ 的概率制备 $\ket{0}$ 或 $\ket{1}$”的\textbf{统计混合}。对后一种制备方法，测量 $X$ 得到 $+1$ 的概率是
\[
  \frac12\abs{\braket{+}{0}}^2+
  \frac12\abs{\braket{+}{1}}^2=\frac12.
\]
而对纯态 $\ket{+}$，概率是 $1$。叠加的概率幅先相加再取模平方；随机混合则按制备概率加权测量概率。一般混合态不能用单个 ket 完整描述，通常用密度算符处理；本文先以纯态为主。

\section{基变换：同一个态可以有不同的坐标}
\subsection{从内积推导变换公式}
设 $\{\ket{e_n}\}$ 和 $\{\ket{f_\alpha}\}$ 是同一有限维空间的两组正交归一基。定义
\begin{equation}
  S_{n\alpha}=\braket{e_n}{f_\alpha},\qquad
  \ket{f_\alpha}=\sum_n\ket{e_n}S_{n\alpha}.
\end{equation}
$S$ 的第 $\alpha$ 列就是新基向量 $\ket{f_\alpha}$ 在旧基中的坐标。由正交归一性可得
\begin{equation}
  (S^\dagger S)_{\alpha\beta}
  =\sum_n\braket{f_\alpha}{e_n}\braket{e_n}{f_\beta}
  =\delta_{\alpha\beta},\qquad S^\dagger S=SS^\dagger=I.
\end{equation}
满足这两个等式的矩阵称为\textbf{酉矩阵}。酉变换保持内积和长度，是复内积空间中正交变换的对应物。

分别记态的旧、新坐标为 $c_n=\braket{e_n}{\psi}$、$d_\alpha=\braket{f_\alpha}{\psi}$，则
\begin{equation}\label{eq:changebasis}
  d_\alpha=\sum_n S_{n\alpha}^*c_n,
  \qquad \boxed{d=S^\dagger c.}
\end{equation}
同样，一个算符在两组基下的矩阵满足
\begin{equation}\label{eq:opchange}
  \boxed{A_{(f)}=S^\dagger A_{(e)}S.}
\end{equation}
因此 $d^\dagger A_{(f)}d=c^\dagger A_{(e)}c$：坐标与矩阵一起变换，物理预言保持不变。

\subsection{二维例子}
从 $\ket{0},\ket{1}$ 换到式\eqref{eq:pm}的 $\ket{+},\ket{-}$，
\begin{equation}
  S=\frac1{\sqrt2}\begin{pmatrix}1&1\\1&-1\end{pmatrix},
  \qquad
  \begin{pmatrix}d_+\\d_-\end{pmatrix}
  =\frac1{\sqrt2}\begin{pmatrix}a+b\\a-b\end{pmatrix}.
\end{equation}
同一个 $X$ 算符在两组基中分别是
\[
  X_{(0,1)}=\begin{pmatrix}0&1\\1&0\end{pmatrix},\qquad
  X_{(+,-)}=\begin{pmatrix}1&0\\0&-1\end{pmatrix}.
\]
非对角元本身不是算符的绝对属性；是否对角取决于所选的基。

\textbf{换基}是改变同一个态的描述；\textbf{主动变换} $\ket{\psi}\mapsto U\ket{\psi}$ 则是在保持描述方式不变时改变状态。两者都可能涉及酉矩阵，但物理含义不同。后文的时间演化属于后者。

\section{位置、动量与连续表示}
\subsection{波函数是态在位置表示中的坐标}
考虑在整条直线上运动、暂不计内部自由度的粒子，其态空间取为 $L^2(\mathbb R)$。形式上引入位置本征态 $\ket{x}$：
\begin{equation}\label{eq:positionbasis}
  \hat x\ket{x}=x\ket{x},\qquad
  \braket{x}{x'}=\delta(x-x'),\qquad
  \int_{\mathbb R}\ket{x}\bra{x}\dd x=\I.
\end{equation}
这里的帽子区分算符 $\hat x$ 与数值 $x$。为避免符号过重，其余算符不一律加帽。

把离散展开中的求和换成积分，得到
\begin{equation}\label{eq:wavefunction}
  \boxed{\psi(x)=\braket{x}{\psi},\qquad
  \ket{\psi}=\int_{\mathbb R}\psi(x)\ket{x}\dd x.}
\end{equation}
这就是\textbf{位置波函数}。它是同一个抽象态的一种表示，而不是在 ket 以外新增的物理对象。

符号 $\delta$ 是 Dirac delta 分布，以积分作用定义：
\[
  \int_{\mathbb R}\delta(x-x_0)f(x)\dd x=f(x_0).
\]
它不是在某一点取“无穷大”的普通函数，不能按普通函数任意相乘。$\ket{x}$ 也不是 $L^2$ 中可归一化的向量，而是\textbf{广义本征态}；式\eqref{eq:positionbasis}是在分布意义下使用的连续基记号。更多的讨论超出了量子力学通常关心的范围。

Born 规则在这里给出\textbf{概率密度}：
\begin{equation}\label{eq:positionBorn}
  \Pr(x\in[a,b])=\int_a^b\abs{\psi(x)}^2\dd x,
  \qquad \int_{\mathbb R}\abs{\psi(x)}^2\dd x=1.
\end{equation}
$\abs{\psi(x)}^2$ 本身不是在一点的概率；连续分布在单个点的概率为零。因为概率无量纲，一维位置波函数的量纲是长度的 $-1/2$ 次方。

\subsection{位置与动量算符怎样作用}
位置算符在位置表示中只是乘以坐标：
\begin{equation}
  (\hat x\psi)(x):=\mel{x}{\hat x}{\psi}=x\psi(x).
\end{equation}
标准的正则动量算符则表示为
\begin{equation}\label{eq:poperator}
  (\hat p\psi)(x)=-\ii\hbar\frac{\dd\psi}{\dd x},
  \qquad \hbar=\frac{h}{2\pi}.
\end{equation}
$h$ 是 Planck 常数，$\hbar$ 具有作用量的量纲。

动量为什么与求导有关？它是\textbf{空间平移的生成元}。若把一个波包主动向右平移 $a$，新波函数为
\[
  (T(a)\psi)(x)=\psi(x-a).
\]
例如原来位于 $x=0$ 的峰值会移到 $x=a$。对足够光滑的函数做无穷小展开，
\begin{equation}
  T(a)\psi=\psi-a\psi'+o(a)
  =\left(\I-\frac{\ii a}{\hbar}\hat p\right)\psi+o(a).
\end{equation}
将动量定义为平移的一阶生成元，就得到式\eqref{eq:poperator}。后面引入算符指数后，这可写成 $T(a)=\exp(-\ii a\hat p/\hbar)$。这给出了动量算符的物理意义；它不是仅凭线性代数推导出的唯一物理定律。

此处必须注意\textbf{定义域}。不是每个 $L^2$ 函数都可以求导，也不是每个 $L^2$ 函数乘上 $x$ 后仍在 $L^2$ 中。因此 $\hat x,\hat p$ 只在适当的稠密子空间上作用。例如整条直线上可取
\[
  D(\hat x)=\{\psi\in L^2:x\psi\in L^2\},\qquad
  D(\hat p)=H^1(\mathbb R),
\]
其中 $H^1$ 指函数及其弱导数都平方可积。暂不熟悉弱导数时，可以先在光滑且快速衰减的波函数上验证公式。

对这样的函数，分部积分给出
\begin{align}
  \braket{\phi}{\hat p\psi}
  &=-\ii\hbar\int_{\mathbb R}\phi^*\psi'\dd x\\
  &=-\ii\hbar[\phi^*\psi]_{-\infty}^{+\infty}
    +\ii\hbar\int_{\mathbb R}(\phi')^*\psi\dd x
  =\braket{\hat p\phi}{\psi}.
\end{align}
边界项消失是所选函数条件的结果。这个计算先验证了对称性；严格的自伴性还要求定义域与伴随的定义域相同。换成有限区间时，边界条件是算符定义的一部分，不能只看 $-\ii\hbar\dd/\dd x$ 的形式就断言自伴。

\subsection{动量表示与 Fourier 变换}
动量本征态满足 $\hat p\ket{p}=p\ket{p}$。其位置波函数 $\braket{x}{p}$ 因而满足
\[
  -\ii\hbar\frac{\dd}{\dd x}\braket{x}{p}=p\braket{x}{p}.
\]
取惯用的相位和归一化约定，
\begin{equation}\label{eq:xpkernel}
  \braket{x}{p}=\frac1{\sqrt{2\pi\hbar}}\ee^{\ii px/\hbar},
  \qquad\braket{p}{p'}=\delta(p-p').
\end{equation}
这些平面波同样不是归一化的 $L^2$ 态；它们是展开波包所用的广义本征态。

动量波函数定义为 $\widetilde\psi(p)=\braket{p}{\psi}$。插入位置表示的单位算符，便得到
\begin{align}
  \widetilde\psi(p)
  &=\int_{\mathbb R}\braket{p}{x}\braket{x}{\psi}\dd x
  =\frac1{\sqrt{2\pi\hbar}}\int_{\mathbb R}
    \ee^{-\ii px/\hbar}\psi(x)\dd x,\label{eq:Fourier}\\
  \psi(x)
  &=\frac1{\sqrt{2\pi\hbar}}\int_{\mathbb R}
    \ee^{\ii px/\hbar}\widetilde\psi(p)\dd p.
\end{align}
\textbf{Fourier 变换正是位置表示与动量表示之间的基变换。}它保持内积，特别是
\[
  \int\abs{\psi(x)}^2\dd x
  =\int\abs{\widetilde\psi(p)}^2\dd p=1.
\]
所以 $\abs{\widetilde\psi(p)}^2$ 是动量的概率密度。一般 $L^2$ 函数的 Fourier 变换按 $L^2$ 极限定义；无需假定上面的积分对每一点都绝对收敛。

在动量表示中，两个算符的角色交换为
\begin{equation}
  \hat p\longleftrightarrow p,
  \qquad \hat x\longleftrightarrow\ii\hbar\frac{\partial}{\partial p}.
\end{equation}
例如对式\eqref{eq:Fourier}求 $p$ 导数，就能验证第二式。算符的物理含义不变，变化的是它在不同表示下的作用形式。

\section{对易关系与不确定性}
\subsection{位置与动量为什么不对易}
两个算符的\textbf{对易子}定义为
\begin{equation}
  \comm{A}{B}=AB-BA.
\end{equation}
在足够光滑且快速衰减的 $\psi$ 上，逐步计算
\begin{align}
  (\hat x\hat p\psi)(x)&=-\ii\hbar x\psi'(x),\\
  (\hat p\hat x\psi)(x)&=-\ii\hbar\frac{\dd}{\dd x}[x\psi(x)]
  =-\ii\hbar\psi(x)-\ii\hbar x\psi'(x).
\end{align}
相减得到正则对易关系
\begin{equation}\label{eq:CCR}
  \boxed{\comm{\hat x}{\hat p}=\ii\hbar\I.}
\end{equation}
它表示两个算符的作用次序不能随便交换。在无限维空间中，这个等式首先是在共同的合适定义域上成立；并不意味着两个乘积可以作用于每个态。

一个有用的检验是：有限维矩阵总满足 $\operatorname{tr}(AB-BA)=0$，而 $\operatorname{tr}(\ii\hbar I)=\ii\hbar N\neq0$。因此有限维矩阵不可能\textbf{精确}满足式\eqref{eq:CCR}；真正的连续位置和动量需要无限维空间。

在有限维中，两个自伴算符对易当且仅当存在一组共同的正交归一本征基。其原因是：若 $A\ket{a}=a\ket{a}$ 且 $[A,B]=0$，则
\[
  A(B\ket{a})=B(A\ket{a})=aB\ket{a},
\]
所以 $B$ 保持 $A$ 的每个本征子空间；再在各子空间内对角化 $B$ 即可。\textbf{不对易不等于没有任何共同本征态}，但意味着不能找到覆盖整个空间的共同本征基。对于 $\hat x,\hat p$，式\eqref{eq:CCR}还排除了共同定义域中的共同本征态。一般无界算符的同时可测性需要更精确的谱投影条件，本文不展开。

\subsection{从 Cauchy--Schwarz 不等式得到不确定关系}
设 $A,B$ 自伴，态已归一化且相关方差与乘积均有定义。记
\[
  \delta A=A-\avg{A}\I,\qquad
  \delta B=B-\avg{B}\I,\qquad
  \ket{u}=\delta A\ket{\psi},\quad
  \ket{v}=\delta B\ket{\psi}.
\]
根据 Cauchy--Schwarz 不等式，
\begin{equation}
  (\Delta A)^2(\Delta B)^2
  =\braket{u}{u}\braket{v}{v}
  \geq\abs{\braket{u}{v}}^2
  \geq\bigl(\operatorname{Im}\braket{u}{v}\bigr)^2.
\end{equation}
而
\[
  2\ii\operatorname{Im}\braket{u}{v}
  =\avg{\delta A\delta B-\delta B\delta A}
  =\avg{[A,B]}.
\]
因此
\begin{equation}\label{eq:uncertainty}
  \boxed{\Delta A\,\Delta B\geq\frac12\abs{\avg{[A,B]}}.}
\end{equation}
代入式\eqref{eq:CCR}即得
\begin{equation}
  \boxed{\Delta x\,\Delta p\geq\frac{\hbar}{2}.}
\end{equation}
这约束的是\textbf{同一种制备状态}下位置分布和动量分布的宽度。可以取许多同样制备的粒子，一部分测位置，另一部分测动量，再计算两组统计分布的标准差；无需先后对同一个粒子测量来定义这两个量。

因此这里的不确定性不是仪器读数误差的别名，也不单是在说第一次测量干扰了第二次测量。它首先是态空间、Born 规则和非对易结构共同给出的限制。对于一般 $A,B$，即使不对易，也可能有某个态使 $\avg{[A,B]}=0$；这时式\eqref{eq:uncertainty}本身不给出正的下界。

\section{动力学假设：薛定谔方程}
\subsection{哈密顿算符决定无测量时的演化}
到目前为止，我们知道给定一个态时怎样计算测量概率，但还不知道这个态以后会怎样变化。

\important{动力学假设}
对于封闭系统，在不进行测量的演化过程中，归一化态满足
\begin{equation}\label{eq:Schrodinger}
  \boxed{\ii\hbar\frac{\dd}{\dd t}\ket{\psi(t)}
  =H(t)\ket{\psi(t)}.}
\end{equation}
$H(t)$ 是系统的自伴\textbf{哈密顿算符}，代表能量，同时生成时间演化。态随时间的变化由 $H$ 与初态共同决定。


在一个\textbf{不随时间变化}的正交归一基中，令 $c_n(t)=\braket{n}{\psi(t)}$，便得到
\begin{equation}\label{eq:Schmatrix}
  \ii\hbar\dot c_m(t)=\sum_n H_{mn}(t)c_n(t).
\end{equation}
这就是一个线性微分方程组。若使用随时间变化的基，对展开式求导时还会出现基向量的导数，不能直接套用式\eqref{eq:Schmatrix}。

\subsection{位置表示中的薛定谔方程}
对于质量为 $m$、在实标量势 $V(x,t)$ 中运动的一维非相对论粒子，常用的哈密顿算符为
\begin{equation}\label{eq:Hparticle}
  H(t)=\frac{\hat p^2}{2m}+V(\hat x,t).
\end{equation}
其中 $V(\hat x,t)$ 在位置表示中就是乘以函数 $V(x,t)$。这个模型把动能与势能的物理信息放进 $H$；复杂系统中如何构造 $H$ 本身也是需要物理判断的问题。

令 $\psi(x,t)=\braket{x}{\psi(t)}$，将动量的微分表示代入，得到
\begin{equation}\label{eq:Schwave}
  \boxed{\ii\hbar\frac{\partial\psi}{\partial t}
  =\left[-\frac{\hbar^2}{2m}\frac{\partial^2}{\partial x^2}
  +V(x,t)\right]\psi.}
\end{equation}
常见的这个偏微分方程，正是抽象方程\eqref{eq:Schrodinger}在位置表示中的写法。

求解问题还需要初态 $\psi(x,t_0)$、空间区域以及合适的边界条件。

\subsection{概率为什么守恒}
由薛定谔方程及 $H^\dagger=H$，
\[
  \frac{\dd}{\dd t}\ket{\psi}=-\frac{\ii}{\hbar}H\ket{\psi},
  \qquad
  \frac{\dd}{\dd t}\bra{\psi}=\frac{\ii}{\hbar}\bra{\psi}H.
\]
所以
\begin{equation}
  \frac{\dd}{\dd t}\braket{\psi}{\psi}
  =\frac{\ii}{\hbar}\mel{\psi}{H}{\psi}
  -\frac{\ii}{\hbar}\mel{\psi}{H}{\psi}=0.
\end{equation}
初态一旦归一化，之后就保持归一化。自伴性同时保证了能量测量的实数结果和这种与概率解释相容的演化。

对于式\eqref{eq:Schwave}，更局部的形式是连续性方程
\begin{equation}
  \frac{\partial\rho}{\partial t}+\frac{\partial j}{\partial x}=0,
  \qquad \rho=\abs{\psi}^2,\qquad
  j=\frac{\hbar}{m}\operatorname{Im}\left(\psi^*\frac{\partial\psi}{\partial x}\right).
\end{equation}
它由薛定谔方程乘上 $\psi^*$、共轭方程乘上 $\psi$ 后相减得到。于是区间内概率的变化率为 $j(a,t)-j(b,t)$，即流入与流出的差；在全空间合适的边界条件下，总概率守恒。

\section{矩阵指数与不含时的时间演化}
\subsection{从数的指数到矩阵的指数}
标量方程 $\dot y=ay$ 的解是 $y(t)=\ee^{at}y(0)$。要解式\eqref{eq:Schmatrix}，自然需要矩阵指数。对任意有限维方阵 $M$，定义
\begin{equation}\label{eq:matrixexp}
  \boxed{\ee^M:=\sum_{k=0}^{\infty}\frac{M^k}{k!}
  =I+M+\frac{M^2}{2!}+\cdots.}
\end{equation}
它\textbf{不是把每个矩阵元分别取指数}。矩阵乘方使用通常的矩阵乘法，而 $M^0=I$。选一个次乘性的矩阵范数，有
\[
  \sum_{k=0}^{\infty}\frac{\norm{M^k}}{k!}
  \leq\sum_{k=0}^{\infty}\frac{\norm{M}^k}{k!}
  =\ee^{\norm{M}}<\infty,
\]
所以定义对每个有限维矩阵都收敛，并可在有界时间区间逐项求导。于是
\begin{equation}\label{eq:expderiv}
  \frac{\dd}{\dd t}\ee^{tM}=M\ee^{tM}=\ee^{tM}M,
  \qquad \ee^{0M}=I.
\end{equation}
这直接证明了 $c(t)=\ee^{tM}c(0)$ 满足 $\dot c=Mc$。

若 $M=SDS^{-1}$，则 $M^k=SD^kS^{-1}$，从而
\begin{equation}
  \ee^M=S\ee^D S^{-1}.
\end{equation}
当 $D=\operatorname{diag}(\lambda_1,\ldots,\lambda_N)$ 时，$\ee^D=\operatorname{diag}(\ee^{\lambda_1},\ldots,\ee^{\lambda_N})$。因此在本征基中取指数最简单，但定义\eqref{eq:matrixexp}并不要求矩阵可对角化。

两个常用性质为
\begin{equation}
  (\ee^M)^\dagger=\ee^{M^\dagger},\qquad
  (\ee^M)^{-1}=\ee^{-M}.
\end{equation}
但一般\textbf{不能}把 $\ee^{A+B}$ 拆成 $\ee^A\ee^B$。展开到二阶就能看出差别：前者含 $(AB+BA)/2$，后者含 $AB$。若 $[A,B]=0$，则可以这样拆分。

\subsection{演化算符及其酉性}
现在假设 $H$ 不显含时间。类比式\eqref{eq:expderiv}，定义
\begin{equation}\label{eq:propagator}
  \boxed{U(t,t_0)=\exp\left[-\frac{\ii}{\hbar}H(t-t_0)\right],
  \qquad\ket{\psi(t)}=U(t,t_0)\ket{\psi(t_0)}.}
\end{equation}
指数中的 $H(t-t_0)/\hbar$ 无量纲。它满足
\[
  \ii\hbar\frac{\partial U}{\partial t}=HU,\qquad
  U(t_0,t_0)=\I,
\]
所以式\eqref{eq:propagator}同时满足运动方程与初始条件。

由于 $H^\dagger=H$，
\begin{equation}
  U^\dagger(t,t_0)=\exp\left[\frac{\ii}{\hbar}H(t-t_0)\right],
  \qquad U^\dagger U=UU^\dagger=\I.
\end{equation}
这称为\textbf{酉演化}。它保持所有态之间的内积，特别是保持范数，因此保持总概率。它还满足
\begin{equation}
  U(t_2,t_1)U(t_1,t_0)=U(t_2,t_0),\qquad
  U(t,t_0)^{-1}=U(t_0,t).
\end{equation}
这些式子分别表达演化可以分段接续，以及封闭系统的这种演化可逆。

\important{从矩阵到无界算符时的界限}
对有界算符，幂级数定义可照搬。但位置、动量及许多哈密顿算符是\textbf{无界算符}：任意态未必能被 $H$ 反复作用，因此不能无条件把式\eqref{eq:matrixexp}作用于任意波函数。对自伴的无界 $H$，式\eqref{eq:propagator}应通过\textbf{谱分解}定义。这样得到的 $U$ 对所有 Hilbert 空间中的态都是良定义的酉算符；若初态属于 $D(H)$，演化才同时以通常的范数可微方式满足薛定谔方程。入门计算中的安全做法，是先求能量本征展开，再给各能量分量乘相位。

\subsection{能量本征态、定态与一般解}
若 $H$ 有完备的离散本征基，
\begin{equation}
  H\ket{n}=E_n\ket{n},\qquad
  H=\sum_n E_n\ket{n}\bra{n},
\end{equation}
那么
\begin{equation}\label{eq:Uspectral}
  U(t,t_0)=\sum_n\ee^{-\ii E_n(t-t_0)/\hbar}\ket{n}\bra{n}.
\end{equation}
对初态 $\ket{\psi(t_0)}=\sum_n c_n\ket{n}$，便有
\begin{equation}\label{eq:generalsolution}
  \boxed{\ket{\psi(t)}=\sum_n c_n
  \ee^{-\ii E_n(t-t_0)/\hbar}\ket{n}.}
\end{equation}
每个能量分量的模不变，因此能量测量的概率不变；不同能量分量之间的\textbf{相对相位}却会变化。这正是很多时间依赖现象的来源。

若初态属于同一个能量本征子空间，整个态只乘上同一个相位，
\[
  \ket{\psi(t)}=\ee^{-\ii E(t-t_0)/\hbar}\ket{\psi(t_0)}.
\]
它代表的射线不变，因此所有不显含时间的可观测量的概率分布不变，称为\textbf{定态}。定态不要求 ket 的数值表达完全不随时间变化。

在位置表示中，令 $\phi_n(x)=\braket{x}{n}$，能量本征方程就是\textbf{定态薛定谔方程}
\begin{equation}
  \left[-\frac{\hbar^2}{2m}\frac{\dd^2}{\dd x^2}+V(x)\right]\phi_n(x)
  =E_n\phi_n(x).
\end{equation}
一般含时解则为
\[
  \psi(x,t)=\sum_n c_n\ee^{-\ii E_n(t-t_0)/\hbar}\phi_n(x).
\]
不要把求一个能量本征态与求任意初态的时间演化混为一件事：前者建立展开基，后者还要使用初态系数 $c_n=\braket{n}{\psi(t_0)}$。

若存在连续谱，相应的求和需要加上对连续谱的积分；不能假定每个哈密顿算符都有完备的可归一化能量本征向量。统一的谱定理写法是
\begin{equation}
  U(t,t_0)=\int_{\mathbb R}\ee^{-\ii E(t-t_0)/\hbar}\dd P_H(E),
\end{equation}
其中 $P_H$ 是 $H$ 的谱投影测度。无需在这里掌握测度论细节；这表示对每个能量部分乘上对应相位，离散情形正好退化为式\eqref{eq:Uspectral}。

\subsection{含时哈密顿算符为什么需要额外小心}
若 $H=H(t)$，在适当的正则性条件下仍可定义演化算符，但通常不能写成
\[
  U(t,t_0)\overset{\text{一般不成立}}{=}
  \exp\left[-\frac{\ii}{\hbar}\int_{t_0}^t H(s)\dd s\right].
\]
原因是不同时间的哈密顿算符可能不对易，先后施加的作用顺序不能丢掉。在有限维等常用情形中，正式写法是\textbf{时间排序指数}
\[
  U(t,t_0)=\mathcal T\exp\left[-\frac{\ii}{\hbar}\int_{t_0}^t H(s)\dd s\right],
\]
$\mathcal T$ 表示在级数展开的乘积中把较晚时刻的算符放在左边。若所有不同时刻的 $H$ 两两对易，就可省去时间排序。本文主要使用不含时情形。

\section{如何计算一个具体量子系统}
\subsection{一个完整的二能级演化例子}
仍用 $\ket{0},\ket{1}$ 作为正交归一基，并取
\begin{equation}
  H=\frac{\hbar\Omega}{2}X
  =\frac{\hbar\Omega}{2}
  \begin{pmatrix}0&1\\1&0\end{pmatrix},\qquad \Omega>0.
\end{equation}
这里 $X$ 是式\eqref{eq:ZX}的算符，$\Omega$ 的量纲为角频率。这只是抽象二能级模型，虽然他确实代表了一类真实系统。设 $t_0=0$，初态为 $\ket{\psi(0)}=\ket{0}$。

\textbf{第一步：求演化算符。}因为 $X^2=\I$，其偶次幂为 $\I$，奇次幂为 $X$。将指数级数按奇偶次分开，得到
\begin{align}
  U(t)&=\ee^{-\ii\Omega tX/2}\\
  &=\left[\sum_{k=0}^{\infty}\frac{(-1)^k(\Omega t/2)^{2k}}{(2k)!}\right]\I
  -\ii\left[\sum_{k=0}^{\infty}\frac{(-1)^k(\Omega t/2)^{2k+1}}{(2k+1)!}\right]X\\
  &=\cos\frac{\Omega t}{2}\,\I
  -\ii\sin\frac{\Omega t}{2}\,X.
\end{align}
\textbf{第二步：作用于初态。}因为 $X\ket{0}=\ket{1}$，
\begin{equation}\label{eq:twoleveltime}
  \ket{\psi(t)}=\cos\frac{\Omega t}{2}\ket{0}
  -\ii\sin\frac{\Omega t}{2}\ket{1}.
\end{equation}
\textbf{第三步：指定测量，再计算概率。}若测量 $Z$，结果 $+1,-1$ 对应的概率分别为
\begin{equation}
  p_Z(+1;t)=\cos^2\frac{\Omega t}{2},\qquad
  p_Z(-1;t)=\sin^2\frac{\Omega t}{2},\qquad
  \avg{Z}(t)=\cos(\Omega t).
\end{equation}
在 $t=\pi/\Omega$ 时，态为 $-\ii\ket{1}$，与 $\ket{1}$ 是同一个物理纯态。

也可以从能量本征基看同一过程：$\ket{+},\ket{-}$ 的能量为 $E_\pm=\pm\hbar\Omega/2$，而
\begin{equation}
  \ket{0}=\frac{\ket{+}+\ket{-}}{\sqrt2},\qquad
  \ket{\psi(t)}=
  \frac{\ee^{-\ii\Omega t/2}\ket{+}+\ee^{\ii\Omega t/2}\ket{-}}{\sqrt2}.
\end{equation}
展开回 $\ket{0},\ket{1}$ 就得到式\eqref{eq:twoleveltime}。能量测量的两个概率始终是 $1/2$，但 $Z$ 测量的概率在变化。\textbf{不含时哈密顿算符保证能量分布不变，并不保证所有物理量的分布不变。}


\subsection{期望的演化，Ehrenfest 定理}
在相关算符和导数均有定义的条件下，对 $\avg{A}(t)=\mel{\psi(t)}{A(t)}{\psi(t)}$ 求导，使用乘积法则与薛定谔方程，得
\begin{align}
  \frac{\dd}{\dd t}\avg{A}
  &=\frac{\ii}{\hbar}\avg{HA}
    +\left\langle\frac{\partial A}{\partial t}\right\rangle
    -\frac{\ii}{\hbar}\avg{AH}\\
  &=\boxed{\frac{1}{\ii\hbar}\avg{[A,H]}
    +\left\langle\frac{\partial A}{\partial t}\right\rangle.}\label{eq:Ehrenfest}
\end{align}
因此，若 $A$ 本身不显含时间且 $[A,H]=0$，则 $\avg{A}$ 守恒。在有限维中，这时 $A$ 的各个谱投影也与演化算符对易，所以整个测量分布都不变。对不含时的 $H$ 取 $A=H$，便得到能量守恒。

对式\eqref{eq:Hparticle}，利用 $[x,p^2]=[x,p]p+p[x,p]=2\ii\hbar p$ 以及 $[p,V(x)]=-\ii\hbar V'(x)$，得到
\begin{equation}
  \frac{\dd\avg{x}}{\dd t}=\frac{\avg{p}}{m},\qquad
  \frac{\dd\avg{p}}{\dd t}=-\avg{V'(x)}.
\end{equation}
这称为 Ehrenfest 关系，说明量子动力学与经典力学如何相接。但一般 $\avg{V'(x)}\neq V'(\avg{x})$，所以平均位置并不总是严格沿着经典轨道运动。

\Needspace{6\baselineskip}
\section{直积与复合体系的态空间}
\important{复合体系假设}
对于两个可区分的子系统 $A,B$，总态空间为张量积
\begin{equation}
  \HH_{AB}=\HH_A\otimes\HH_B.
\end{equation}
若两个空间分别有基 $\{\ket{i}_A\}$、$\{\ket{j}_B\}$，则总空间以 $\ket{i}_A\otimes\ket{j}_B$ 为基，简写为 $\ket{ij}$。有限维时，总维数是两者维数的乘积。

一般纯态写成
\begin{equation}
  \ket{\Psi}=\sum_{i,j}c_{ij}\ket{i}_A\otimes\ket{j}_B,
  \qquad \sum_{i,j}\abs{c_{ij}}^2=1.
\end{equation}
若系数可分解为 $c_{ij}=a_i b_j$，态可以写成 $\ket{\psi}_A\otimes\ket{\phi}_B$，称为乘积态。不能这样分解的纯态称为\textbf{纠缠态}。例如
\begin{equation}
  \ket{\Psi}=\frac{\ket{00}+\ket{11}}{\sqrt2}
\end{equation}
不能写成两个单体系 ket 的张量积，可通过比较四个展开系数验证。

只作用于子系统 $A$ 的可观测量 $O_A$，在总空间中表示为 $O_A\otimes\I_B$。若分别测两个子系统，事件对应投影 $P_a\otimes Q_b$，则
\begin{equation}
  p(a,b)=\mel{\Psi}{P_a\otimes Q_b}{\Psi}.
\end{equation}
仍然是同一条 Born 规则。纠缠态中，单个子系统一般不能由自己的一个 ket 描述；要继续处理它，需要密度算符和偏迹。本文到此只保留复合空间的基本结构。


\end{document}
